\subsection*{Why a finite source block cannot reduce the unsigned prime norm}

The arithmetic bound in the preceding tail estimates is not merely a
poor choice of positive weight.  There is an exact recurrence obstruction
to replacing it by a small norm on the source complement.

\begin{proposition}[Prime norm survives every finite-dimensional removal]
\label{prop:v119-prime-essential}
Fix $\lambda>1$ and let $\mathsf T_{\rm pr}$ be the positive-coefficient
sum of truncated prime translations in
\eqref{eq:v116-weighted-prime-general}.  For every finite-rank
orthogonal projection $P$ on $L^2(0,L)$, $Q=I-P$, one has
\begin{equation}\label{eq:v119-prime-compression}
 \|Q\mathsf T_{\rm pr}Q\|=\|\mathsf T_{\rm pr}\|.
\end{equation}
For every compact operator $C$ on this space,
\begin{equation}\label{eq:v119-prime-compact}
 \|Q(\mathsf T_{\rm pr}-C)Q\|\ge\|\mathsf T_{\rm pr}\|.
\end{equation}
In particular the essential norm of $\mathsf T_{\rm pr}$ equals its
ordinary norm.  The projection need not be a Fourier projection and
may remove any prescribed finite source block.
\end{proposition}
\begin{proof}
For each positive integer $k$, let
$M_k f(x)=e^{2\pi ikx/L}f(x)$.  Simultaneous recurrence of the finitely
many phases gives integers $k_j\to\infty$ such that
\[
 e^{2\pi ik_j\log m/L}\longrightarrow1
 \quad\hbox{for every active }m.
\]
For completeness, partition the unit cube in the finitely many phase
coordinates into $H^d$ cubes and compare $H^d+1$ successive multiples.
Two lie in one cube, giving a positive return with coordinate errors
at most $1/H$.  If such returns have an unbounded subsequence, use it;
otherwise a fixed positive return has all coordinate errors zero,
and its increasing multiples give the required sequence.  No rational
independence of prime logarithms is assumed.

Conjugating a truncated translation by $M_k$ only multiplies it by
its corresponding phase.  Since each translation has norm at most one,
\begin{equation}\label{eq:v119-prime-recurrence}
 \|M_{k_j}^*\mathsf T_{\rm pr}M_{k_j}-\mathsf T_{\rm pr}\|
 \le2\sum_m w_m|e^{2\pi ik_j\log m/L}-1|\longrightarrow0.
\end{equation}
For every fixed $f\in L^2(0,L)$, $M_{k_j}f\rightharpoonup0$:
each scalar product is a Fourier coefficient of an $L^1$ function.
Thus $PM_{k_j}f\to0$, $PM_{k_j}\mathsf T_{\rm pr}f\to0$, and
$CM_{k_j}f\to0$ strongly.  Combining these statements with
\eqref{eq:v119-prime-recurrence} gives
\[
 Q(\mathsf T_{\rm pr}-C)Q M_{k_j}f
       -M_{k_j}\mathsf T_{\rm pr}f\longrightarrow0.
\]
For unit $f$, the input has norm at most one, so taking norms and then
the supremum over $f$ proves \eqref{eq:v119-prime-compact}.
Putting $C=0$ and using compression contractivity proves
\eqref{eq:v119-prime-compression}.  Taking the infimum over compact
$C$, with $P=0$, proves the essential-norm assertion.
\end{proof}

\begin{proposition}[The sharp growing-window scale of the unsigned prime norm]
\label{prop:v119-prime-scale}
With the actual weights $w_m=\Lambda(m)/\sqrt m$ and $L=2\log\lambda$,
\begin{equation}\label{eq:v119-prime-scale}
 \|\mathsf T_{\rm pr}\|=(1+o(1))\lambda
       \qquad(\lambda\longrightarrow\infty).
\end{equation}
Consequently \eqref{eq:v119-prime-scale} also holds for
$\|Q_\lambda\mathsf T_{\rm pr}Q_\lambda\|$ for any family of
finite-rank orthogonal projections $I-Q_\lambda$, with no restriction on their
finite ranks or Fourier cutoffs.
\end{proposition}
\begin{proof}
Write
\[
 a(x)=e^{-x/2},\quad b(x)=e^{-(L-x)/2},\quad\phi(x)=a(x)+b(x),
\]
and
\[
 \Psi(X)=\sum_{1<m<X}\Lambda(m),\quad
 S(X)=\sum_{1<m<X}\frac{\Lambda(m)}m.
\]
The strict cutoff agrees with the truncated translations almost
everywhere.  Direct evaluation gives
\begin{equation}\label{eq:v119-prime-weight-ratio}
 \frac{(\mathsf T_{\rm pr}\phi)(x)}{\phi(x)}
 =\frac{a(x)\{S(e^{L-x})+\Psi(e^x)\}
       +b(x)\{\Psi(e^{L-x})+S(e^x)\}}{a(x)+b(x)}.
\end{equation}
The exact identity
\[
 a(x)e^x+b(x)e^{L-x}=\lambda\{a(x)+b(x)\}
\]
identifies the leading constant.  The classical prime number
theorem~\cite{NISTPrimeAsymptotics}, equivalently $\Psi(X)\sim X$,
implies that for each $\delta>0$ there is $C_\delta<\infty$ with
\[
 |\Psi(X)-X|\le\delta X+C_\delta\qquad(X\ge1).
\]
The Chebyshev bound and partial summation give
$0\le S(X)\le C(1+\log X)$, with an absolute $C$.
Therefore, uniformly for $0<x<L$,
\[
 (1-\delta)\lambda-C_\delta
 \le\frac{\mathsf T_{\rm pr}\phi}{\phi}
 \le(1+\delta)\lambda+C_\delta+C(1+L).
\]
At each fixed window $\phi$ is bounded and bounded away from zero.
The upper estimate and the weighted Schur inequality
\eqref{eq:v116-weighted-prime-general} bound the operator norm from
above.  Integrating the lower estimate against $\phi^2$ gives the
same lower bound for its Rayleigh quotient and hence for the norm.
Divide by $\lambda$, let $\lambda\to\infty$, and then let
$\delta\downarrow0$. This proves \eqref{eq:v119-prime-scale}.
The projection assertion follows from
Proposition~\ref{prop:v119-prime-essential} at each window.
\end{proof}

These statements concern the unsigned prime operator, not the sign of
the full Weil form.  They make precise a limitation of the particular
place-by-place bound \eqref{eq:v115-sharp-tail}. Even replacing its
prime constant by the exact tail norm cannot certify positivity unless
\begin{equation}\label{eq:v119-unsigned-cutoff}
 \log\frac{N+1}{L}>\|\mathsf T_{\rm pr}\|
       =(1+o(1))\lambda.
\end{equation}
Here we refer to the whole-space bound, whose remaining subtracted
terms are nonnegative. Thus this method still requires an exponential
Fourier cutoff in $\lambda$, rather than the polynomial source cutoff.
This is a limitation of that sufficient scalar bound, not a necessary
cutoff for actual Weil positivity or for a frequency-weighted argument.

The physical pole operator has rank at most two.  Equation
\eqref{eq:v119-prime-compact} shows that combining it with the prime
operator does not make their ordinary operator norm smaller than
$\|\mathsf T_{\rm pr}\|$, even after a finite source removal.
The recurrence frequencies in the proof may be arbitrarily large;
there the archimedean logarithmic energy also grows.  No negative
direction for the complete Weil operator is inferred.  A scalable G2
estimate must control the arithmetic term together with that energy
or retain more precise signed correlations, rather than assert a small
unweighted norm on the complement.  The required signed estimate
remains open.
